\clearpage\thispagestyle{empty}
\pdfbookmark[0]{Quellen zum Brief an Bode}{bode1810sep-sources}
\begin{center}\textbf{Quellen zum Brief an Bode}\end{center}
\noindent Diese KI-Übertragung ins heutige Deutsch folgt dem Brief an Johann Elert Bode vom 2. September 1810 in \href{https://gdz.sub.uni-goettingen.de/id/PPN236007467}{Gauß’ \textit{Werken}, Bd.~VI} (Göttingen, 1874), S.~319. Bode druckte den Auszug auf S.~\href{https://api.digitale-sammlungen.de/iiif/image/v2/bsb10538342_00245/full/1600,/0/default.jpg}{237} des \textit{Astronomischen Jahrbuchs für 1813} (Berlin, 1810) nach der Ephemeride auf S.~\href{https://api.digitale-sammlungen.de/iiif/image/v2/bsb10538342_00244/full/1600,/0/default.jpg}{236}--237 ab. Auf der folgenden S.~\href{https://api.digitale-sammlungen.de/iiif/image/v2/bsb10538342_00246/full/1600,/0/default.jpg}{238} steht eine gesonderte, von Bode gezeichnete Notiz. Das digitalisierte Exemplar stammt aus der Bayerischen Staatsbibliothek.
\par\smallskip\noindent Die eckige Klammer zur Identifizierung der Ephemeride stammt aus der Werkausgabe. Im Jahrbuch steht ein Punkt nach der Bemerkung, der Empfänger wisse bereits von den Plänen der Regierung. In der Werkausgabe folgt der Betrag von 200000 Franken an dieser Stelle ohne Satzzeichen. KI-Übertragung und KI-Satz.
\clearpage
